\chapter{Mathematical Background} \label{ch:background}

This chapter collects the mathematical material needed later for randomized
least-squares methods. The main ingredients come from three directions:
classical least-squares theory, random projections and subspace embeddings, and
perturbation questions related to conditioning. Although all spaces in this thesis
are finite-dimensional, it is often convenient to use the language of linear maps,
ranges, and subspaces rather than only matrix coordinates.

For an overdetermined system \(A \in \R^{n \times d}\) with \(n \ge d\), the least-squares problem
\[
	\min_{x \in \R^d} \|Ax-b\|_2
\]
asks for the vector whose image under \(A\) is closest to \(b\) in Euclidean norm.
Geometrically, this means projecting \(b\) onto the range of \(A\). If \(A\) has
full column rank, the minimizer is unique. In exact arithmetic, it may be
characterized by the normal equations or, more stably, by an orthogonal
factorization.

The randomized viewpoint replaces the original problem by a smaller one. A
sketching matrix \(S \in \R^{s \times n}\), with \(s\) much smaller than \(n\), is
chosen so that the geometry of the range of \(A\) is approximately preserved after
compression. One then solves
\[
	\min_{x \in \R^d} \|SAx-Sb\|_2.
\]
The quality of this replacement depends on how well the sketch preserves norms
and angles on the relevant subspace. This is where subspace embeddings and the
Johnson--Lindenstrauss phenomenon enter the picture. The sketch should act
almost isometrically on the subspace that carries the least-squares geometry. The
formal definitions and embedding guarantees are given in \Cref{ch:randls}.

Two sketch families are especially important in this thesis. The first is the
Gaussian sketch, whose entries are independent random variables with suitable
scaling. The second is the class of structured sketches, especially transforms based
on randomized Hadamard matrices. Gaussian sketches are the cleanest theoretical
baseline. Structured sketches are more attractive computationally because they can
exploit fast transforms.

Conditioning plays a central role throughout. Even when sketching gives a good
approximation in principle, the numerical behavior of the reduced problem depends
strongly on the conditioning of \(SA\), that is, on the stability of the compressed
operator on the relevant subspace. This becomes even more important in finite
precision, since rounding errors may be harmless for a well-conditioned problem
and severe for an ill-conditioned one. The relation between sketching quality and
conditioning is discussed in \Cref{prop:sketched_cond}.

The later mixed-precision discussion therefore rests on a simple principle.
Randomized approximation error and floating-point error should not be studied in
isolation. A useful theory must account for both at the same time.
